\Article{237}
\textsc{Theorem.} \textit{If a form $F$ is transformable into the product of two forms $f,f'$, and the form $f'$ contains the form $f''$, then $F$ is also transformable into the product of the forms $f,f''$.}

\textit{Proof.} Keep all the notation of Article 235 for the forms $F,f,f'$; let the form $f''$ be $=(a'',b'',c'')$, and let $f'$ pass into $f''$ by the substitution $\alpha.\ \beta,\gamma,\delta$. Then it is easy to see that $F$ passes into $ff''$ by the substitution
\[
\begin{array}{llll}
\alpha p+\gamma p',&\beta p+\delta p',&\alpha p''+\gamma p''',&\beta p''+\delta p'''\\
\alpha q+\gamma q',&\beta q+\delta q',&\alpha q''+\gamma q''',&\beta q''+\delta q'''
\end{array}
\qquad\text{\textit{Q. E. D.}}
\]

For brevity, put the coefficients
\[
\begin{gathered}\alpha p+\gamma p',\quad\beta p+\delta p'\ \text{etc.}\\=\mathfrak{P},\mathfrak{P}',\mathfrak{P}'',\mathfrak{P}''';\quad\mathfrak{Q},\mathfrak{Q}',\mathfrak{Q}'',\mathfrak{Q}'''\end{gathered}
\]
\SourceLeaf{257}
and the number $\alpha\delta-\beta\gamma=e$; from equations $\Omega$ of Article 235 we readily verify that
\[
\begin{aligned}
\mathfrak{P}\mathfrak{Q}'-\mathfrak{Q}\mathfrak{P}'&=an'e\\
\mathfrak{P}\mathfrak{Q}'''-\mathfrak{Q}\mathfrak{P}'''-\mathfrak{P}'\mathfrak{Q}''+\mathfrak{Q}'\mathfrak{P}''&=2bn'e\\
\mathfrak{P}''\mathfrak{Q}'''-\mathfrak{Q}''\mathfrak{P}'''&=cn'e\\
\mathfrak{P}\mathfrak{Q}''-\mathfrak{Q}\mathfrak{P}''&=\alpha\alpha a'n+2\alpha\gamma b'n+\gamma\gamma c'n=a''n\\
\mathfrak{P}\mathfrak{Q}'''-\mathfrak{Q}\mathfrak{P}'''+\mathfrak{P}'\mathfrak{Q}''-\mathfrak{Q}'\mathfrak{P}''&=2b''n\\
\mathfrak{P}'\mathfrak{Q}'''-\mathfrak{Q}'\mathfrak{P}'''&=c''n\\
\mathfrak{Q}'\mathfrak{Q}''-\mathfrak{Q}\mathfrak{Q}'''&=Ann'e\\
\mathfrak{P}\mathfrak{Q}'''+\mathfrak{Q}\mathfrak{P}'''-\mathfrak{P}'\mathfrak{Q}''-\mathfrak{Q}'\mathfrak{P}''&=2Bnn'e\\
\mathfrak{P}'\mathfrak{P}''-\mathfrak{P}\mathfrak{P}'''&=Cnn'e
\end{aligned}
\]
Now, denoting the determinant of the form $f''$ by $d''$, we have that $e$ is a square root of $\frac{d''}{d'}$, positive or negative according as the form $f'$ contains the form $f''$ properly or improperly. So $n'e$ is a square root of $\frac{d''}{D}$. Thus the preceding nine equations are clearly entirely analogous to equations $\Omega$ of Article 235, and the form $f$ is taken in the transformation of the form $F$ into $ff''$ in the same way as in the transformation of the form $F$ into $ff'$; but the form $f''$ is taken in the former either in the same way as $f'$ in the latter, or in the opposite way, according as $f'$ contains $f''$ properly or improperly.

\Needspace{8\baselineskip}
\Article{238}
\textsc{Theorem.} \textit{If a form $F$ is contained in a form $F'$ and is transformable into the product of the forms $f,f'$, then the form $F'$ is also transformable into the same product.}

\textit{Proof.} Keeping the same notation as above for the forms $F,f,f'$, and supposing that the form $F'$ passes into $F$ by the substitution $\alpha,\beta,\gamma,\delta$, we readily see that $F'$, by the substitution
\[
\begin{array}{llll}
\alpha p+\beta q,&\alpha p'+\beta q',&\alpha p''+\beta q'',&\alpha p'''+\beta q'''\\
\gamma p+\delta q,&\gamma p'+\delta q',&\gamma p''+\delta q'',&\gamma p'''+\delta q'''
\end{array}
\]
becomes the same as $F$ does by the substitution $p,p',p'',p''';\ q,q',q'',q'''$, and therefore $F'$ passes into $ff'$ by that substitution. \textit{Q. E. D.}

Further, a calculation like the one in the preceding article readily confirms that $F'$ is transformable into $ff'$ in the same way as $F$, when $F'$ contains $F$ properly; but when $F$ is improperly contained in $F'$, the transformations of the form $F$ into $ff'$ and of the form $F'$ into $ff'$ are opposite with respect to each of the forms $f,f'$: namely,
\SourceLeaf{258}
whichever of these forms enters one transformation directly is taken inversely in the other.

Combining this theorem with the theorem of the preceding article gives the following more general one: \textit{If a form $F$ is transformable into the product $ff'$, and the forms $f,f'$ contain the forms $g,g'$, respectively, while the form $F$ is contained in the form $G$, then $G$ is transformable into the product $gg'$.} For by the theorem of this article, $G$ is transformable into $ff'$, then by the theorem of the preceding article into $fg'$, and by the same theorem also into $gg'$. Further, if all three forms $f,f',G$ contain the forms $g,g',F$ properly, then $G$ is clearly transformable into $gg'$ in the same way with respect to the forms $g,g'$ as $F$ into $ff'$ with respect to the forms $f,f'$; the same happens if all three containments are improper. Finally, it is just as easy to determine \textit{how} $G$ is transformable into $gg'$ if one of these containments differs from the other two.

If the forms $F,f,f'$ are equivalent to the forms $G,g,g'$, respectively, the latter have the same determinants as the former, and for the forms $f,f'$, the numbers $m,m'$ are the same for the forms $g,g'$ (Article 161). From this, by the fourth conclusion of Article 235, we deduce without difficulty that in this case $G$ is \textit{composed} of $g,g'$ if $F$ is composed of $f,f'$. Also, the form $g$ enters the former composition in the same way as $f$ enters the latter, when $F$ is equivalent to $G$ in the same way as $f$ to $g$, and conversely; similarly, $g'$ is taken in the former composition in the same or the opposite way as $f'$ in the latter, according as the equivalence of the forms $f',g'$ is of the same or a different kind as the equivalence of the forms $F,G$.

\Needspace{8\baselineskip}
\Article{239}
\textsc{Theorem.} \textit{If a form $F$ is composed of the forms $f,f'$, then every other form transformable into the product $ff'$ in the same way as $F$ contains $F$ properly.}

\textit{Proof.} Keeping all the notation of Article 235 for $F,f,f'$, equations $\Omega$ hold here too. Suppose the form $F'=(A',B',C')$, whose determinant is $=D'$, passes into the product $ff'$ by the substitution $\mathfrak{p},\mathfrak{p}',\mathfrak{p}'',\mathfrak{p}''';\ \mathfrak{q},\mathfrak{q}',\mathfrak{q}'',\mathfrak{q}'''$, and denote the numbers
\[
\begin{gathered}\mathfrak{p}\mathfrak{q}'-\mathfrak{q}\mathfrak{p}',\quad
\mathfrak{p}\mathfrak{q}''-\mathfrak{q}\mathfrak{p}'',\quad
\mathfrak{p}\mathfrak{q}'''-\mathfrak{q}\mathfrak{p}'''.\\
\mathfrak{p}'\mathfrak{q}''-\mathfrak{q}'\mathfrak{p}'',\quad
\mathfrak{p}'\mathfrak{q}'''-\mathfrak{q}'\mathfrak{p}''',\quad
\mathfrak{p}''\mathfrak{q}'''-\mathfrak{q}''\mathfrak{p}'''\end{gathered}
\]
by
\[
P',Q',R',S',T',U'
\]
\SourceLeaf{259}
respectively. Then we have nine equations entirely analogous to $\Omega$, namely
\[
\begin{gathered}
P'=a\mathfrak{n}',\quad R'-S'=2b\mathfrak{n}',\quad U'=c\mathfrak{n}'\\
Q'=a'\mathfrak{n},\quad R'+S'=2b'\mathfrak{n},\quad T'=c'\mathfrak{n}\\
\mathfrak{q}'\mathfrak{q}''-\mathfrak{q}\mathfrak{q}'''=A'\mathfrak{n}\mathfrak{n}',\\
\mathfrak{p}\mathfrak{q}'''+\mathfrak{q}\mathfrak{p}'''-\mathfrak{p}'\mathfrak{q}''-\mathfrak{q}'\mathfrak{p}''=2B'\mathfrak{n}\mathfrak{n}',\\
\mathfrak{p}'\mathfrak{p}''-\mathfrak{p}\mathfrak{p}'''=C'\mathfrak{n}\mathfrak{n}'
\end{gathered}
\]
which we will denote by $\Omega'$. Here the quantities $\mathfrak{n},\mathfrak{n}'$ are square roots of $\frac{d}{D'},\frac{d'}{D'}$, taken with the same signs as $n,n'$, respectively; so if the positive square root of $\frac{D}{D'}$ (which is an integer) is put $=k$, then $\mathfrak{n}=kn$, $\mathfrak{n}'=kn'$. From this and the first six equations in $\Omega$ and $\Omega'$, we clearly have
\[
\begin{gathered}
P'=kP,\quad Q'=kQ,\quad R'=kR\\
S'=kS,\quad T'=kT,\quad U'=kU
\end{gathered}
\]
So by the lemma of Article 234, four integers $\alpha,\beta,\gamma,\delta$ can be determined such that
\[
\begin{aligned}
\alpha p+\beta q&=\mathfrak{p}.&\gamma p+\delta q&=\mathfrak{q}\\
\alpha p'+\beta q'&=\mathfrak{p}',&\gamma p'+\delta q'&=\mathfrak{q}'\ \text{etc.}
\end{aligned}
\]
and
\[
\alpha\delta-\beta\gamma=k
\]
Substituting these values of $\mathfrak{p},\mathfrak{q},\mathfrak{p}',\mathfrak{q}'$, etc. into the last three equations $\Omega'$, we readily verify, with the help of the equations $\mathfrak{n}=kn,\ \mathfrak{n}'=kn'$ and the last three of $\Omega$, that
\[
\begin{aligned}
A'\alpha\alpha+2B'\alpha\gamma+C'\gamma\gamma&=A\\
A'\alpha\beta+B'(\alpha\delta+\beta\gamma)+C'\gamma\delta&=B\\
A'\beta\beta+2B'\beta\delta+C'\delta\delta&=C
\end{aligned}
\]
Thus the form $F'$ passes by the substitution $\alpha,\beta,\gamma,\delta$ (which is proper, since $\alpha\delta-\beta\gamma=k$ is positive) into $F$, \textit{i.e.}, it contains the form $F$ properly. \textit{Q. E. D.}

So if $F'$ is also composed of the forms $f,f'$ (in the same way as $F$ is composed of them), the forms $F,F'$ have the same determinant, and are therefore properly equivalent. More generally, if a form $G$ is composed of the forms $g,g'$ in the same way as
\SourceLeaf{260}
$F$ of $f,f'$, respectively, and the forms $g,g'$ are properly equivalent to $f,f'$, then the forms $F,G$ are properly equivalent.

Since the case in which both forms enter the composition directly is the simplest, and the others are easily reduced to it, we will consider only this case below. Thus whenever a form is simply said to be composed of two others, we must always understand that it is properly composed of each of them\HistoricalNote{*)}{Similarly, in the composition of ratios (which has a close analogy with the composition of forms), it is usually understood that the ratios to be composed are to be taken directly, unless the opposite is stated.}. The same restriction will apply whenever a form is said to be transformable into the product of two others.
